\section{Instantiation}\label{sec:inst}

We describe a concrete instantiation of the black box construction in pairing
based groups using algebraic PRFs and tailored sigma protocols, yielding a
transparent setup.

\subsection{Building blocks}

\begin{primitivebox}{BBS+ signatures}{fig:bbs}

  \begin{description}

    \item[$(\sk,\pk) \sample \BBSgen(\secparam, L)$ :] On input the security
      parameter and vector length $L$, do the following:

      \begin{enumerate}

        \item Sample $z \sample \Zp$ and generators
          $\gen_0,\gen_1,\ldots,\gen_{L+1} \sample \group{1}$, $\ggen \sample
          \group{2}$.

        \item Return $\sk = z$, $\pk = (\gen_0,\ldots,\gen_{L+1},\ggen,\ggen_Z
          = \ggen^z)$.

      \end{enumerate}

    \item[$\signature \sample \BBSsign(\sk,(M_1,\ldots,M_L))$ :] On input
      secret key $\sk = z$ and message $(M_1,\ldots,M_L) \in \group{1}^L$, do:

      \begin{enumerate}

        \item Sample $e,r \sample \Zp$.

        \item Return $\signature = (e,r,E)$ with $E = (\gen_0 \cdot
          \prod_{i=1}^{L} M_i \cdot \gen_{L+1}^{r})^{1/(z+e)}$.

      \end{enumerate}

    \item[$\{0,1\} \leftarrow \BBSverify(\pk,(M_1,\ldots,M_L),\signature)$ :]
      On input public key $\pk$, message $(M_1,\ldots,M_L)$, and signature
      $\signature = (e,r,E)$, do:

      \begin{enumerate}

        \item Return $e(E,\ggen_Z\cdot\ggen^{e}) \checkcheck e(\gen_0 \cdot
          \prod_{i=1}^{L} M_i \cdot \gen_{L+1}^{r}, \ggen)$.

      \end{enumerate}

  \end{description}

\end{primitivebox}

This variant of BBS+ signs the group elements $M_i = \gen_i^{m_i}$ directly,
rather than the scalars $m_i$ themselves. Because computing $\prod_i M_i$
requires no knowledge of the discrete logarithms $m_i$, a signer can produce a
valid signature on committed values without ever learning them; this is exactly
what lets the registration authority certify a user's keys below without ever
seeing them~\cite{SCN:AuSusMu06,EC:TesZhu23b}.

\begin{primitivebox}{Pedersen commitments}{fig:pedersen}

  \begin{description}

    \item[$\ck \sample \ComGen(\secparam)$ :] On input the security parameter,
      sample $\gen,\hgen \sample \group{1}$ and return $\ck = (\gen,\hgen)$.

    \item[$\com \sample \ComCommit(\ck,m,\openinfo)$ :] On input commitment key
      $\ck = (\gen,\hgen)$, message $m \in \Zp$, and opening $\openinfo \in
      \Zp$, return $\com = \gen^{m} \cdot \hgen^{\openinfo}$.

    \item[$m/\bot \leftarrow \ComOpen(\ck,\com,m,\openinfo)$ :] On input $\ck$,
      commitment $\com$, message $m$, and opening $\openinfo$, return $m$ if
      $\com = \gen^{m}\cdot\hgen^{\openinfo}$, and $\bot$ otherwise.

  \end{description}

\end{primitivebox}

Pedersen commitments are perfectly hiding, so nobody, not even the registration
authority or an auditor, learns anything about a committed tracing key beyond
what is explicitly disclosed.  Commitments are group elements, so users can
prove statements using cheap sigma protocols over discrete logarithms rather
than requiring general purpose zero knowledge proofs.

\subsection{Our transaction identification scheme}\label{sec:algebraic}

\begin{protocolbox}{Transaction identification protocol}{fig:txid}

  \begin{description}

    \item[$(\crs,\state) \sample \IDgen(\secparam)$ :] On input the security
      parameter, do the following:

      \begin{enumerate}

        \item Sample bilinear groups $(\group{1},\group{2},\group{T},p,e)$ of
          prime order $p$ admitting an efficient pairing $e: \group{1} \times
          \group{2} \to \group{T}$.

        \item Sample generators $\gen_0,\gen_1,\gen_2,\gen_3 \sample \group{1}$
          and $\ggen,\ggen_1,\ldots,\ggen_\ell \sample \group{2}$, and fix a
          hash function $\HE : \{0,1\}^* \to \group{1}$, modelled as a random
          oracle.

        \item The registration authority samples $(\sk_{RA},\pk_{RA}) \sample
          \BBSgen(\secparam,2)$.

        \item Sample $\ck \sample \ComGen(\secparam)$ and set $\ck \leftarrow
          (\gen_2,\gen_3)$, reusing the generators of Step 2.

        \item Set $\Max \leftarrow 2^\ell$, $\FE(\utk,\epoch) :=
          \HE(\epoch)^{\utk}$, and $\FT(\ek,\gamma) := e(\ek,
          \prod_{i=1}^{\ell} \ggen_i^{\gamma_i})$ over the bit decomposition
          $(\gamma_1,\ldots,\gamma_\ell)$ of $\gamma$.

        \item Set $\crs \leftarrow (\group{1},\group{2},\group{T},e,
          \gen_0,\gen_1,\gen_2,\gen_3,\ggen,\ggen_1,\ldots,\ggen_\ell,\HE,
          \pk_{RA},\ck,\Max)$ and $\state \leftarrow \varnothing$, the empty
          ledger.

        \item Return $(\crs,\state)$.

      \end{enumerate}

    \item[$(\sk,\rc) \sample \IDregister(\crs,\state)$ :] Interactively between
      the user and the registration authority, do:

      \begin{enumerate}

        \item The registration authority sends a fresh nonce $\rand$ to the
          user.

        \item The user samples $\sk,\utk \sample \Zp$, sets $S \leftarrow
          \gen_1^{\sk}$, $T \leftarrow \gen_2^{\utk}$, computes $\soksig
          \sample \soksign(\crs_0,\inst_0,\witness_0,\rand)$ for $\inst_0 =
          (\gen_1,\gen_2,S,T)$ and $\witness_0 = (\sk,\utk)$, proving $S =
          \gen_1^{\sk} \wedge T = \gen_2^{\utk}$, and sends $(S,T,\soksig)$ to
          the registration authority.

        \item If $\soksig$ does not verify, or the user has registered before,
          the registration authority exits.

        \item Else, the registration authority computes $\signature \sample
          \BBSsign(\sk_{RA},(S,T))$ and records $(S,T)$ as registered.

        \item Return $(\sk,\utk)$ as the user's secret key, and $\rc \leftarrow
          (S,T,\signature)$ as their registration credential.

      \end{enumerate}

    \item[$\ek \sample \IDgrant(\crs,\sk,\utk,\rc,\epoch)$ :] Interactively
      between the user and an auditor, on a fresh nonce $\rand$ from the
      auditor, do:

      \begin{enumerate}

        \item Parse $\rc = (S,T,\signature)$.

        \item The user computes $\ek \leftarrow \FE(\utk,\epoch) =
          \HE(\epoch)^{\utk}$.

        \item The user samples $\alpha \sample \Zp$, computes $(A,B) =
          (\gen_2^{\alpha},\HE(\epoch)^{\alpha})$, sets $c =
          \hash(\gen_2,\epoch,A,B)$ and $\eta_1 = \alpha + c\cdot\utk$, forming
          $\pi \leftarrow (A,B,\eta_1)$, proving $\rel_1$: $\ek =
          \HE(\epoch)^{\utk} \wedge T = \gen_2^{\utk}$.

        \item The user computes $\psi \sample
          \soksign(\crs_0,\inst_0,\witness_0,\pi\|\rand)$, re-authenticating
          $(S,T)$ on the fresh nonce, and sends $(S,T,\signature,\ek,\pi,\psi)$
          to the auditor.

        \item The auditor accepts and outputs $\ek$ iff $\psi$ verifies,
          $\BBSverify(\pk_{RA},(S,T),\signature) = 1$, and $\pi$ verifies for
          $\rel_1$.

      \end{enumerate}

    \item[$\tx \sample \IDsubmit(\crs,\sk,\utk,\rc,\epoch,\payload)$ :] On
      input $(\crs,\sk,\utk,\rc,\epoch,\payload)$ from the user, do:

      \begin{enumerate}

        \item Parse $\rc = (S,T,\signature)$ and compute $\ek \leftarrow
          \FE(\utk,\epoch)$.

        \item Sample a fresh, unused counter $0 \leq \gamma < \Max$ and compute
          $\Tag_\gamma \leftarrow \FT(\ek,\gamma)$.

        \item Sample $\openinfo \sample \Zp$ and compute $\com \sample
          \ComCommit(\ck,\utk,\openinfo)$.

        \item Compute $\pi_1 \sample \NIZKprove(\crs_3,\inst_3,\witness_3)$,
          proving relation $\rel_3$, following~\Cref{sec:tag-computation}.

        \item Compute $\psi_2 \sample \soksign(\crs_2,\inst_2,\witness_2,\msg)$
          for $\msg = \payload\|\Tag_\gamma\|\com\|\pi_1$, proving relation
          $\rel_2$, following~\Cref{sec:tx-auth}.

        \item Set $\tx \leftarrow (\payload,\Tag_\gamma,\com,\pi_1,\psi_2)$ and
          submit $\tx$ to the ledger.

        \item The ledger appends $\tx$ to $\state$ iff $\Tag_\gamma$ has not
          appeared in a previously accepted transaction, $\pi_1$ verifies for
          $\rel_3$, and $\psi_2$ verifies for $\rel_2$ on $\msg$.

        \item Return $\tx$.

      \end{enumerate}

    \item[$\{\payload\} \leftarrow \IDidentify(\crs,\ek,\epoch)$ :] On input
      $(\crs,\ek,\epoch)$, do:

      \begin{enumerate}

        \item For each $0 \leq \gamma < \Max$, compute $\Tag_\gamma \leftarrow
          \FT(\ek,\gamma)$ and retrieve $\state[\Tag_\gamma]$ if present.

        \item Return the set of retrieved payloads.

      \end{enumerate}

  \end{description}

\end{protocolbox}

The protocol instantiation is described in full in~\Cref{fig:txid}. The
remaining algorithms follow~\Cref{sec:txidconstruction} directly; we highlight
below only the design choices specific to this instantiation.

The registration authority's key pair is for BBS+, chosen because it signs
group elements directly (\Cref{fig:bbs}), letting the registration authority
certify a user's keys without ever seeing them in the clear. This is essential
since the registration authority must not learn a user's tracing key. During
registration, the user accordingly sends only the group elements $S =
\gen_1^{\sk}$ and $T = \gen_2^{\utk}$, together with a proof of knowledge of
the underlying secrets, rather than the secrets themselves.

The commitment scheme is Pedersen, chosen because it is perfectly hiding, and
because its algebraic structure lets a user prove statements about a
commitment using cheap sigma protocols rather than general purpose zero
knowledge proofs. The commitment key reuses the pair of generators
$(\gen_2,\gen_3)$ that appear in the registration authority's public key,
purely for compactness; randomness used at signing time and at commitment time
remains independently sampled at each use.

Deriving the epoch tracing key deterministically from the long term tracing
key and the epoch identifier is what confines an auditor's tracing power to a
single epoch, since a key derived for one epoch is useless for any other. Each
transaction tag is computed by evaluating a pairing based PRF at the user's
epoch key and a fresh counter, so it looks random to anyone without that key.
The tag carries two separate proofs, that it was correctly derived from a
properly committed tracing key (\Cref{sec:tag-computation}), and that the
committed key was signed by the registration authority
(\Cref{sec:tx-auth}).

\subsubsection*{Verifiable tag computation.}\label{sec:tag-computation} The
user first computes Pedersen commitment $\com =
\gen_2^{\utk}\cdot\gen_3^{\openinfo}$ and $\Tag = \FT(\ek, \gamma)$, where $
\FT(\ek, \gamma)$ is defined as $e\!\left(\ek,\, \prod_{i=1}^{\ell}
\ggen_i^{\gamma_i}\right)$.  Under the DDH assumption, and the Bilinear
Decisional Diffie Hellman (BDDH) assumption, $\FT$ can be shown to be a PRF, in
the random oracle model.

Relation $\rel_3$ is defined for $\inst_3 = (\HE, \{\ggen_i\}_{i\in[1,\ell]},
\gen_2, \gen_3, \com, \Max, \epoch, \Tag)$; $\witness_3 = (\utk, \openinfo,
\ek, \{\gamma_i\}_{i\in[1,\ell]})$; and $(\varnothing, \crs_3)\leftarrow
\NIZKgen(\rel_3)$ with $\crs_3$ describing a hash function
$\hash:\{0,1\}^*\rightarrow \Zp$, and $\rel_3(\inst_3, \witness_3) =1$ if and
only if the following holds: \begin{equation*} \com =
  \gen_2^{\utk}\cdot\gen_3^{\openinfo} \ \wedge \ \ek = \HE(\epoch)^{\utk}\
\wedge \ \Tag_{\gamma} = e(\ek,\prod_{i=1}^{\ell} \ggen_i^{\gamma_i}) \ \wedge
\ \forall i \in[1, \ell]: \gamma_i \in \{0,1\}.  \end{equation*}

Proving $\rel_3$ directly would require an argument that proves quadratic
equations over pairings, which would be expensive in practice.

Instead, the user randomly samples $r, s \leftarrow \Zp$, and with them
computes auxiliary values $\Gamma = \ggen^r\prod_{i=1}^{\ell}
\ggen_i^{\gamma_i}, U = H(\epoch)^s, V = \ek^{s}$. The user then proves the
equivalent relation $\widehat{\rel}_3$ defined for $\widehat{\inst}_3
=(\inst_3, \ggen, \Gamma, U, V)$; $\widehat{\witness}_3 = (\witness_3, r, s)$;
and $\widehat{\crs}_3 = \crs_3$, that together verify: \begin{gather*} \com =
  \gen_2^{\utk}\cdot\gen_3^{\openinfo} \quad \wedge \quad \Gamma =
  \ggen^r\prod_{i=1}^{\ell} \ggen_i^{\gamma_i} \quad \wedge \quad \forall i \in
  [1,\ell]: \gamma_i \in \{0,1\} \\ \wedge \quad U = \HE(\epoch)^s \quad \wedge
  \quad V = U^{\utk} \quad \wedge \quad e(V, \Gamma) = \Tag^s \cdot e(V,
\ggen)^r.  \end{gather*}

We prove this in two parts: a sigma protocol for the equalities of
$\widehat{\rel}_3$ involving $\widehat{\witness}_3$, and a separate range
argument that $(\gamma_1,\ldots,\gamma_\ell)$ are bits.

\begin{protocolbox}{Tag correctness sigma protocol}{txid-tag-sigma}

  \begin{description}

    \item[$\pi \sample \NIZKprove(\crs_3, \widehat{\inst}_3,
      \widehat{\witness}_3)$ :] On input $\widehat{\inst}_3 = (\com, \Gamma, U,
      V, \Tag)$ and $\widehat{\witness}_3 = (\utk, \openinfo, r, s)$, do:

      \begin{enumerate}

        \item Sample $\alpha, \beta, \eta, \nu \sample \Zp$ and $\delta_i
          \sample \Zp$ for each $i \in [1,\ell]$.

        \item Compute $A = \gen_2^{\alpha}\cdot\gen_3^{\beta}$, $\Delta =
          \ggen^{\eta}\prod_{i=1}^{\ell} \ggen_i^{\delta_i}$, $B =
          \HE(\epoch)^{\nu}$, $D = U^{\alpha}$, and $F = \Tag^{\nu} \cdot e(V,
          \ggen)^{\eta}$.

        \item Compute challenge $c \leftarrow \hash(\widehat{\inst}_3, A,
          \Delta, B, D, F)$.

        \item Compute $\theta = \alpha + c\cdot\utk$, $\omega = \beta +
          c\cdot\openinfo$, $\upsilon = \eta + c\cdot r$, $\xi_i = \delta_i +
          c\cdot\gamma_i$ for each $i \in [1,\ell]$, and $\chi = \nu +
          c\cdot s$.

        \item Return $\pi = (A, \Delta, B, D, F, \theta, \omega, \{\xi_i\},
          \upsilon, \chi)$.

      \end{enumerate}

    \item[$\{0,1\} \leftarrow \NIZKverify(\crs_3, \widehat{\inst}_3, \pi)$ :]
      On input $\widehat{\inst}_3$ and proof $\pi$, compute challenge $c
      \leftarrow \hash(\widehat{\inst}_3, A, \Delta, B, D, F)$ and return the
      result of the following checks:
      %
      \begin{gather*} \gen_2^{\theta}\cdot\gen_3^{\omega} \checkcheck A \cdot
        \com^c, \quad \ggen^{\upsilon}\prod_{i=1}^{\ell} \ggen_i^{\xi_i}
        \checkcheck \Delta \cdot \Gamma^c, \quad \HE(\epoch)^{\chi} \checkcheck
        B \cdot U^c, \\ U^{\chi} \checkcheck D \cdot V^c, \quad \Tag^\chi \cdot
        e(V, \ggen)^\upsilon \checkcheck F \cdot e(V,\Gamma)^c.
      \end{gather*}

  \end{description}

\end{protocolbox}

For each $i\in[\ell]$, bit validity of $\gamma_i$ is established using an inner
product argument~\cite{SP:BBBPWM18}.

  \subsubsection*{Transaction authorisation.}\label{sec:tx-auth} Each
  transaction carries a signature of knowledge that attests it has originated
  from a registered user, following relation $\rel_2$ defined below.

  Since the user is registered using BBS+ signature $\signature = (e, r, E =
  ({\gen_0}\cdot\gen_2^{\sk} \cdot \gen_1^{\utk}\cdot\gen_3^r)^{1/(z+e)})$,
  $\rel_2$ is defined for $\inst_2 = (\pk, \com) = (\gen_0, \gen_1, \gen_2,
  \gen_3, \ggen, \ggen_Z, \com)$, $\witness_2 = (\sk, \utk, \rho, e, r, E)$,
  $\rel_2(\inst_2, \witness_2) = 1$ if and only if the following holds:
  \begin{equation*} \com = \gen_2^{\utk}\cdot\gen_3^{\rho} \ \wedge \ e(E,
    \ggen^e\cdot\ggen_Z) =
  e({\gen_0}\cdot\gen_2^{\sk}\cdot\gen_1^{\utk}\cdot\gen_3^r, \ggen),
\end{equation*} and $\crs_2 \leftarrow \sokgen(\rel_2)$ is set to hash function
$\hash:\{0, 1\}^*\rightarrow \Zp$.

The SoK for $\rel_2$ on message $\msg$ is obtained by rewriting $\rel_2$ as an
equivalent relation $\widehat{\rel}_2$ that avoids proofs over bilinear
equations~\cite{EPRINT:CamDriLeh16,USENIX:HesSinSor23}.  $\widehat{\rel}_2$ is
defined for  $\widehat{\inst}_2 = (\pk, \com, U, V, W)$ defined as follows:
\begin{equation*} U = E^{\zeta/\eta}; \quad V=
({\gen_0}\cdot{\gen_2}^{\sk}\cdot{\gen_1}^{\utk}\cdot{\gen_3}^r)^{ 1/\zeta};
\quad W= U^eV^{\zeta}, \end{equation*} and $\widehat{\witness}_2 = (\sk, \utk,
\rho, e, r', \zeta, \eta)$ such that $\widehat{\rel}_2(\widehat{\inst}_2,
\widehat{\witness}_2) = 1$ if and only if the following holds:
\begin{equation*} \com = \gen_2^{\utk}\cdot\gen_3^{\rho}\ \wedge \ e(U,
\ggen_Z)\cdot{e(W, \ggen)}=1  \ \wedge \ W = U^e\cdot{V^{\zeta}} \ \wedge \
(\gen_0\cdot\com)^{ 1} = V^{\eta}\cdot \gen_1^{\sk}\cdot\gen_3^{r'}.
\end{equation*}

$\widehat{\rel}_2$ can be proven using sigma protocols, and the Fiat Shamir
heuristic transforms the resulting proof into a signature of knowledge.

\begin{protocolbox}{Transaction authorisation signature of knowledge}{txid-tx-sok}

  \begin{description}

    \item[$\soksig \sample \soksign(\crs_2, \widehat{\inst}_2,
      \widehat{\witness}_2, \msg)$ :] On input $\widehat{\inst}_2 = (\pk,
      \com, U, V, W)$, $\widehat{\witness}_2 = (\sk, \utk, \rho, e, r',
      \zeta, \eta)$, and message $\msg$, do:

      \begin{enumerate}

        \item Sample $\alpha, \beta, \delta, \upsilon, \chi, \xi, \theta
          \sample \Zp$.

        \item Compute $A = \gen_2^{\alpha}\cdot\gen_3^{\beta}$, $B =
          U^{\delta}\cdot V^{\upsilon}$, and $D =
          V^{\chi}\cdot\gen_1^{\xi}\cdot \gen_3^{\theta}$.

        \item Compute challenge $c \leftarrow \hash(\widehat{\inst}_2, A, B,
          D, \msg)$.

        \item Compute $\phi_1 = \alpha + c\cdot\utk$, $\phi_2 = \beta +
          c\cdot\rho$, $\phi_3 = \delta + c\cdot e$, $\phi_4 = \upsilon +
          c\cdot\zeta$, $\phi_5 = \chi + c\cdot\eta$, $\phi_6 = \xi +
          c\cdot\sk$, and $\phi_7 = \theta + c\cdot r'$.

        \item Return $\soksig = (U, V, W, A, B, D, \vec{\phi})$, for
          $\vec{\phi} = (\phi_1,\ldots,\phi_7)$.

      \end{enumerate}

    \item[$\{0,1\} \leftarrow \sokverify(\crs_2, \com, \msg, \soksig)$ :] On
      input commitment $\com$, message $\msg$, and SoK $\soksig$, do:

      \begin{enumerate}

        \item If $e(U, \ggen_Z)\cdot e(W, \ggen) \neq 1$, return $0$.

        \item Compute challenge $c \leftarrow \hash(\widehat{\inst}_2, A, B,
          D, \msg)$.

        \item Return the result of the following checks:
          %
          \begin{gather*} A\cdot\com^c \checkcheck
            \gen_2^{\phi_1}\cdot\gen_3^{\phi_2}, \quad B\cdot W^c \checkcheck
            U^{\phi_3}V^{\phi_4}, \\ D\cdot(\gen_0\cdot\com)^c \checkcheck
            V^{\phi_5}\cdot \gen_1^{\phi_6}\cdot\gen_3^{\phi_7}.
          \end{gather*}

      \end{enumerate}

  \end{description}

\end{protocolbox}

\begin{assumptionbox}{Decisional Diffie-Hellman (DDH)}{ddh-txid}

  The DDH assumption holds in $\group{1}$ if no PPT adversary can distinguish
  $(\gen^x, \gen^y, \gen^{xy})$ from $(\gen^x, \gen^y, \gen^z)$ for uniformly
  random $x, y, z \sample \Zp$.

\end{assumptionbox}

\begin{assumptionbox}{Bilinear Decisional Diffie-Hellman (BDDH)}{bddh}

  The BDDH assumption holds if, given $(\gen, \gen^a, \gen^b, \ggen, \ggen^c)$
  for uniformly random $a, b, c \sample \Zp$, no PPT adversary can distinguish
  $e(\gen,\ggen)^{abc}$ from a uniformly random element of $\group{T}$.

\end{assumptionbox}

\begin{assumptionbox}{$q$-Strong Diffie-Hellman ($q$-SDH)}{qsdh-txid}

  The $q$-SDH assumption holds if, given $(\gen, \gen^{x}, \gen^{x^2}, \dots,
  \gen^{x^q}, \ggen, \ggen^{x})$ for uniformly random $x \sample \Zp$, no PPT
  adversary can output a pair $(c, \gen^{1/(x+c)})$ for any $c \in \Zp
  \setminus \{-x\}$ with non-negligible probability.

\end{assumptionbox}

\begin{theorem}\label{thm:inst} The instantiation presented
in~\Cref{sec:algebraic} securely realises $\idfu{txid}$ in the random oracle
model, under the Decisional Diffie Hellman assumption in $\group{1}$, the
Bilinear Decisional Diffie Hellman assumption, and the $q$-Strong Diffie
Hellman assumption.  \end{theorem}
